\cnabstracttitle

在新能源高比例并网和煤电机组深度调峰背景下，燃煤机组汽轮机长期处于宽负荷、频繁变工况运行状态，其压力、温度、流量等热工参数不仅受设备健康状态影响，还随负荷水平、蒸汽边界及运行调节持续变化。同时，汽轮机系统具有明显的多设备、多参数耦合特征，而现场高质量故障样本数量有限，传统固定阈值及单纯依赖故障样本的数据驱动方法难以兼顾变工况适应性、故障识别精度和诊断解释性。本文以燃煤机组汽轮机热力系统为研究对象，围绕设备健康状态监测、故障知识图谱构建及KG-GCN故障诊断方法开展研究。主要研究内容包括：

针对汽轮机正常运行参数随工况变化而持续迁移、固定设计基准难以直接用于长期状态评价的问题，提出基于GRU的设备级健康状态监测方法。该方法依据设备实际热力边界，将操作变量和上游/外部边界作为输入，以设备本体及接口状态参数作为输出，通过时序模型学习给定运行条件下的健康响应关系。以超高压缸（VHP）为代表性对象，采用2024年5月正常DCS数据训练模型，并以2024年6月数据进行独立测试，同时设置ANN、TCN和iTransformer-small进行对比。结果表明，GRU在最终24个状态参数上的宏平均$R^2$达到0.9511，在VHP本体、一级抽汽及一次冷再接口等主要状态上均能够形成较稳定的动态健康参考，为后续利用实测值与健康预测值之间的偏差表征设备状态变化提供了基础。

针对汽轮机设备关系复杂、故障知识分散于DCS点表、设备结构资料及故障文献中的问题，构建基于运行生产资料的汽轮机系统故障知识图谱。根据汽轮机主通流、再热、抽汽回热及凝汽冷端等实际热力过程，定义设备、参数、异常、故障和处置等实体类型以及相应关系类型；采用大语言模型生成候选三元组，并结合实体类型规则、关系方向、本机设备兼容性、来源证据及Embedding语义对齐完成知识校核和实体标准化。当前审计知识库包含127个标准实体和123条候选关系，其中121条关系进入主知识图谱，2条存在资料差异的候选关系保留待核。由此将分散的设备结构信息与故障机理转化为可用于图卷积计算的结构化领域知识。

针对不同汽轮机故障可能共享局部监测参数以及故障标签数量有限的问题，提出基于知识图谱与图卷积网络的KG-GCN故障诊断方法。该方法将健康状态模型产生的多参数残差按时间窗口映射为知识图谱参数节点的动态特征，并将知识图谱投影得到的参数关系作为图结构输入，通过GCN沿设备和参数关联关系逐层聚合故障特征。训练阶段首先在不使用类别标签的残差窗口上进行随机遮蔽重构预训练，随后固定图卷积编码器并利用带标签故障窗口训练分类层。以VHP为代表性算例，基于真实健康运行残差和独立故障机理构造正常状态、VHP通流部件侵蚀、一级抽汽支路泄漏和VHP汽封泄漏四类半物理诊断场景。固定随机种子实验中，KG-GCN的诊断Accuracy和F1分别达到98.57\%和98.58\%，优于CNN和AE；3组随机种子重复实验的平均Macro-F1达到98.46\%$\pm$0.11\%。在分类阶段仅使用20\%带标签故障窗口时，掩码预训练使Macro-F1由94.48\%提高至96.67\%，表明知识拓扑与残差预训练能够在故障标签有限条件下提高诊断性能和结果稳定性。

综上，本文形成了由“设备健康状态建模—动态状态残差—汽轮机故障知识图谱—图卷积特征聚合—掩码预训练—故障分类”组成的汽轮机系统状态监测与故障诊断流程。本文故障分类实验采用真实健康残差背景与独立机理构造的半物理故障场景，相关结果验证了所提方法在当前代表性算例中的有效性；其面向现场真实故障及其他汽轮机设备的工程泛化能力仍需结合后续故障记录和经实测数据校准的仿真样本进一步验证。

\cnkeywords{燃煤机组；汽轮机系统；状态监测；知识图谱；图卷积网络；故障诊断}
